% Slug: ch35-guard-point
% Source: chapters/ch35.tex:33 (tikzpicture env start line)
% Original label: \label{fig:ch35-guard-point}
% Original caption (LaTeX source):
%   Each guarded operation on the fast path can reach a shared exit stub that rebuilds interpreter state, so no bet is placed without a tested way back.
% Original caption (plain text):
%   Each guarded operation on the fast path can reach a shared exit stub that rebuilds interpreter state, so no bet is placed without a tested way back.
% Required TikZ setup: \usepackage{tikz} + \usetikzlibrary{arrows.meta, positioning, shapes.geometric, fit} (loaded by support/compiler-pseudocode.sty).
% Shared styles used in this picture: `basic block', `cfg edge'.
%   Defined in support/compiler-pseudocode.sty: `basic block', `cfg edge', `back edge', `block label'.
% Usage: % Slug: ch35-guard-point
% Source: chapters/ch35.tex:33 (tikzpicture env start line)
% Original label: \label{fig:ch35-guard-point}
% Original caption (LaTeX source):
%   Each guarded operation on the fast path can reach a shared exit stub that rebuilds interpreter state, so no bet is placed without a tested way back.
% Original caption (plain text):
%   Each guarded operation on the fast path can reach a shared exit stub that rebuilds interpreter state, so no bet is placed without a tested way back.
% Required TikZ setup: \usepackage{tikz} + \usetikzlibrary{arrows.meta, positioning, shapes.geometric, fit} (loaded by support/compiler-pseudocode.sty).
% Shared styles used in this picture: `basic block', `cfg edge'.
%   Defined in support/compiler-pseudocode.sty: `basic block', `cfg edge', `back edge', `block label'.
% Usage: % Slug: ch35-guard-point
% Source: chapters/ch35.tex:33 (tikzpicture env start line)
% Original label: \label{fig:ch35-guard-point}
% Original caption (LaTeX source):
%   Each guarded operation on the fast path can reach a shared exit stub that rebuilds interpreter state, so no bet is placed without a tested way back.
% Original caption (plain text):
%   Each guarded operation on the fast path can reach a shared exit stub that rebuilds interpreter state, so no bet is placed without a tested way back.
% Required TikZ setup: \usepackage{tikz} + \usetikzlibrary{arrows.meta, positioning, shapes.geometric, fit} (loaded by support/compiler-pseudocode.sty).
% Shared styles used in this picture: `basic block', `cfg edge'.
%   Defined in support/compiler-pseudocode.sty: `basic block', `cfg edge', `back edge', `block label'.
% Usage: % Slug: ch35-guard-point
% Source: chapters/ch35.tex:33 (tikzpicture env start line)
% Original label: \label{fig:ch35-guard-point}
% Original caption (LaTeX source):
%   Each guarded operation on the fast path can reach a shared exit stub that rebuilds interpreter state, so no bet is placed without a tested way back.
% Original caption (plain text):
%   Each guarded operation on the fast path can reach a shared exit stub that rebuilds interpreter state, so no bet is placed without a tested way back.
% Required TikZ setup: \usepackage{tikz} + \usetikzlibrary{arrows.meta, positioning, shapes.geometric, fit} (loaded by support/compiler-pseudocode.sty).
% Shared styles used in this picture: `basic block', `cfg edge'.
%   Defined in support/compiler-pseudocode.sty: `basic block', `cfg edge', `back edge', `block label'.
% Usage: \input{graphs-code/ch35-guard-point.tikz.tex} inside a figure environment
%   (wrap with \begin{figure} ... \caption{...} \label{fig:ch35-guard-point} ... \end{figure}).
% Note: tikzpicture body below is verbatim from the source; do not reformat.
\begin{tikzpicture}[node distance=1.0cm and 1.6cm]
\node[basic block] (fast1) {guarded op 1};
\node[basic block, below=of fast1] (fast2) {guarded op 2};
\node[basic block, below=of fast2] (fast3) {guarded op 3};
\node[basic block, right=of fast2] (stub) {exit stub\\ rebuild state};
\node[basic block, right=of stub] (interp) {interpreter\\ resume};
\draw[cfg edge] (fast1) -- (fast2);
\draw[cfg edge] (fast2) -- (fast3);
\draw[cfg edge] (fast1) -- (stub);
\draw[cfg edge] (fast2) -- (stub);
\draw[cfg edge] (fast3) -- (stub);
\draw[cfg edge] (stub) -- (interp);
\end{tikzpicture}
 inside a figure environment
%   (wrap with \begin{figure} ... \caption{...} \label{fig:ch35-guard-point} ... \end{figure}).
% Note: tikzpicture body below is verbatim from the source; do not reformat.
\begin{tikzpicture}[node distance=1.0cm and 1.6cm]
\node[basic block] (fast1) {guarded op 1};
\node[basic block, below=of fast1] (fast2) {guarded op 2};
\node[basic block, below=of fast2] (fast3) {guarded op 3};
\node[basic block, right=of fast2] (stub) {exit stub\\ rebuild state};
\node[basic block, right=of stub] (interp) {interpreter\\ resume};
\draw[cfg edge] (fast1) -- (fast2);
\draw[cfg edge] (fast2) -- (fast3);
\draw[cfg edge] (fast1) -- (stub);
\draw[cfg edge] (fast2) -- (stub);
\draw[cfg edge] (fast3) -- (stub);
\draw[cfg edge] (stub) -- (interp);
\end{tikzpicture}
 inside a figure environment
%   (wrap with \begin{figure} ... \caption{...} \label{fig:ch35-guard-point} ... \end{figure}).
% Note: tikzpicture body below is verbatim from the source; do not reformat.
\begin{tikzpicture}[node distance=1.0cm and 1.6cm]
\node[basic block] (fast1) {guarded op 1};
\node[basic block, below=of fast1] (fast2) {guarded op 2};
\node[basic block, below=of fast2] (fast3) {guarded op 3};
\node[basic block, right=of fast2] (stub) {exit stub\\ rebuild state};
\node[basic block, right=of stub] (interp) {interpreter\\ resume};
\draw[cfg edge] (fast1) -- (fast2);
\draw[cfg edge] (fast2) -- (fast3);
\draw[cfg edge] (fast1) -- (stub);
\draw[cfg edge] (fast2) -- (stub);
\draw[cfg edge] (fast3) -- (stub);
\draw[cfg edge] (stub) -- (interp);
\end{tikzpicture}
 inside a figure environment
%   (wrap with \begin{figure} ... \caption{...} \label{fig:ch35-guard-point} ... \end{figure}).
% Note: tikzpicture body below is verbatim from the source; do not reformat.
\begin{tikzpicture}[node distance=1.0cm and 1.6cm]
\node[basic block] (fast1) {guarded op 1};
\node[basic block, below=of fast1] (fast2) {guarded op 2};
\node[basic block, below=of fast2] (fast3) {guarded op 3};
\node[basic block, right=of fast2] (stub) {exit stub\\ rebuild state};
\node[basic block, right=of stub] (interp) {interpreter\\ resume};
\draw[cfg edge] (fast1) -- (fast2);
\draw[cfg edge] (fast2) -- (fast3);
\draw[cfg edge] (fast1) -- (stub);
\draw[cfg edge] (fast2) -- (stub);
\draw[cfg edge] (fast3) -- (stub);
\draw[cfg edge] (stub) -- (interp);
\end{tikzpicture}
